% !PRINTED none
% [sic] no printed folio number visible on this section-opening page (Zehnter Abschnitt); page contains the Abschnitt title, the \S.\,108 heading, and running body text, but no running head/number is legible in any of the three bands.

\begin{center}
\gesperrt{Zehnter Abschnitt.}

\textbf{Gen\"aherte Berechnung der Wurzeln.}

\rule{3cm}{0.4pt}

\S.\ 108.

\gesperrt{Interpolation. Regula falsi.}
\end{center}

Mit der Abgrenzung der Intervalle, in denen nur je eine Wurzel einer algebraischen Gleichung liegt, ist die M\"oglichkeit gegeben, die reellen Wurzeln mit beliebiger Genauigkeit numerisch zu berechnen, indem man das Intervall mehr und mehr einengt, z. B. fortgesetzt halbirt. Theilt man ein Intervall von der Gr\"osse 1 fortgesetzt in zehn Theile, so liefert jeder neue Schritt eine weitere Decimalstelle.

Wenn einmal erst eine Wurzel von allen \"ubrigen abgesondert ist, so hat man bei der weiteren Einengung des Intervalles immer nur das Vorzeichen der Function $f(x)$ selbst zu ber\"ucksichtigen. Die dazu n\"othigen Rechnungen k\"onnen sehr erleichtert werden durch die Anwendung der Interpolationsformel, die wir im ersten Abschnitt kennen gelernt haben.

Wir haben n\"amlich in \S.\,10, (7) eine Formel hergeleitet, nach der man eine Function $n^{\mathrm{ten}}$ Grades, die wir jetzt mit $\varphi(\xi)$ bezeichnen wollen, bestimmen kann, wenn $n+1$ auf einander folgende Werthe $\varphi(0), \varphi(1) \ldots \varphi(n)$ gegeben sind. Diese Formel ist, wenn
\weq{1}{\begin{aligned}
B_\nu^{(\xi)} &= \frac{\xi(\xi-1)\ldots(\xi-\nu+1)}{1.2.3\ldots\nu}\\
\Delta_\xi &= \varphi(\xi+1) - \varphi(\xi),
\end{aligned}}
%CONT>>
